%\thispagestyle{empty}
%\pagestyle{custom}% Set page style to custom
%\pdfbookmark[0]{Declaration}{declaration} % Bookmark name visible in a PDF viewer
% \begin{center}
% 	\Huge{\textbf{STATEMENT}}
% \end{center}
\pagestyle{empty}% Set page style to empty
\pdfbookmark[0]{Declaration}{declaration}% Bookmark name visible in a PDF viewer
\begin{center}
	\Huge \textbf{Declaration}
	\hrulefill
\end{center}
\begin{flushleft}
	{\large
	\noindent I, \myName, hereby declare that this thesis has not been, and will not be, submitted in whole or in part to another university for the award of any other degree.
	
	The work presented in this thesis is divided into two research projects and has been
	completed in collaboration with Andrea Banfi, Melissa van Beekveld,
	Silvia Ferrario Ravasio, Alexander Karlberg, Matthew Lim, and Daniel Reichelt.
	
	The content of chapter~\ref{chapter:DIS} is taken nearly verbatim from the paper~\cite{vanBeekveld:2026rez}
	currently undergoing peer review:
	
	M. van Beekveld, S. Ferrario Ravasio, A. Karlberg, D. Peake \textit{A generalised-$k_t$ jet algorithm for Deep
	Inelastic Scattering} arXiv:2606.13077
	
	I produced the first implementation of the algorithm and carried out the accompanying
	literature review, and performed the majority of the phenomenological runs on LXPLUS
	using Pythia and Rivet. Later implementations were developed together
	with the other authors, with Melissa preparing the version that was released in the plugin.
	The plots were produced by Silvia from my runs, and the conceptual development was
	shared between all of us. Support and supervision were given throughout by all.
	
	Appendix~\ref{ap:DIS-fastjet} contains relevant parts of the code taken from the
	\href{https://repo.hepforge.org/source/fastjetsvn/browse/contrib/contribs/DISGenkt}{\texttt{DISGenkt}} plugin.
	
	The content of chapter~\ref{ch:hhZjet} is based on a project that will be submitted for publication in due
	course:
	
	D. Peake, A. Banfi, M.A. Lim, D. Reichelt \textit{NNLL resummation in Z+jet events with the ARES method}
	
	The majority of the calculations in this project were performed by me, under the supervision
	of Andrea throughout and with further input from Matthew and Daniel. In particular, I computed
	all NNLL corrections to the observable, and with Andrea derived the initial-state treatment in the 
	\textsc{ARES} formalism. I performed all validations for the work presented here.
	\par}
\end{flushleft}
\vspace*{2cm}
\begin{minipage}{.45\linewidth}
	\begin{flushleft} 
		\textit{\myLocation,} \\
		\textit{\today}
	\end{flushleft}
\end{minipage}
\hfill
\begin{minipage}{.45\linewidth}
	\begin{flushright} 
		\makebox[2.5in]{\hrulefill} \\
		\myName 
	\end{flushright}
\end{minipage}\\ [0.5cm]